\documentclass[aspectratio=169,11pt]{beamer}
\usepackage[T1]{fontenc}
\usepackage[utf8]{inputenc}
\usepackage[brazilian]{babel}
\usepackage{lmodern,booktabs,tikz}
\usetikzlibrary{arrows.meta,positioning}
\definecolor{navy}{HTML}{10233F}
\definecolor{teal}{HTML}{087E8B}
\definecolor{gold}{HTML}{D6A437}
\definecolor{muted}{HTML}{526176}
\definecolor{paper}{HTML}{F3F6FA}
\setbeamercolor{normal text}{fg=navy,bg=white}
\setbeamercolor{structure}{fg=teal}
\setbeamercolor{frametitle}{fg=navy,bg=paper}
\setbeamercolor{block title}{fg=white,bg=navy}
\setbeamercolor{block body}{fg=navy,bg=paper}
\setbeamerfont{frametitle}{size=\Large,series=\bfseries}
\setbeamertemplate{navigation symbols}{}
\setbeamertemplate{itemize items}[circle]
\setbeamersize{text margin left=8mm,text margin right=8mm}
\setbeamertemplate{footline}{\hspace{8mm}\color{muted}\tiny Aristides H. G. da Cruz\quad|\quad Orientador: Gustavo Luís Soares\hfill DyraSQL\quad\insertframenumber\hspace{8mm}\vspace{3mm}}
\setbeamertemplate{frametitle}{\vspace{4mm}\insertframetitle\par\vspace{1mm}{\color{gold}\rule{13mm}{1.2mm}}\vspace{1mm}}
\setbeamerfont{block title}{size=\normalsize,series=\bfseries}
\setbeamerfont{block body}{size=\normalsize}
\newcommand{\source}[1]{\par\vspace{2mm}{\scriptsize\color{muted}Fonte: #1}}
\newcommand{\evidence}[2]{\begin{center}\includegraphics[width=\textwidth,height=.66\textheight,keepaspectratio]{evidencias/#1}\end{center}\source{#2}}
\hypersetup{pdftitle={DyraSQL - Pitch de TCC},pdfauthor={Aristides Henrique Gonçalves da Cruz; Orientador: Gustavo Luís Soares}}
\title{DyraSQL}
\author{Aristides Henrique Gonçalves da Cruz}
\date{13 de outubro de 2026}
\begin{document}
% Pitch: slides 1 a 12. Roteiro de gravação em roteiro_pitch.md (6 minutos).
\begin{frame}[plain]
\vspace{4mm}
{\color{teal}\small PUC MINAS\quad/\quad SISTEMAS DE INFORMAÇÃO}
\vspace{7mm}

{\fontsize{38}{42}\selectfont\bfseries DyraSQL}\par
\vspace{3mm}
{\LARGE Roteamento dinâmico de consultas SQL}\par
\vspace{2mm}
{\large Baseado em metadados de tabelas Apache Iceberg}
\vspace{5mm}

{\color{gold}\rule{21mm}{1.5mm}}\par
\vspace{4mm}
Aristides Henrique Gonçalves da Cruz\par
{\small Orientador: Gustavo Luís Soares}\par
\vspace{2mm}
{\small\color{muted}Trabalho de Conclusão de Curso\quad|\quad Outubro de 2026}
\end{frame}

\begin{frame}{Contexto e motivação}
\begin{columns}[T,onlytextwidth]
\column{.48\textwidth}
\begin{block}{Dados compartilhados}
Data lakes separam armazenamento e processamento. O mesmo conjunto de dados pode ser consultado por ambientes de capacidades distintas.
\end{block}
\column{.48\textwidth}
\begin{block}{Demandas diferentes}
Consultas variam em volume, junções, agregações e filtros. Uma configuração fixa pode superdimensionar consultas simples ou limitar as complexas.
\end{block}
\end{columns}
\vspace{4mm}
\begin{block}{Pergunta de pesquisa}
Como usar informações disponíveis antes da execução para escolher entre três faixas de capacidade do Trino?
\end{block}
\source{Artigo do TCC; Strausz, Pardon e Giurgiu (2025).}
\end{frame}

\begin{frame}{Objetivo e contribuição}
\begin{block}{Objetivo geral}
Desenvolver um mecanismo de decisão que combine metadados e características SQL para selecionar o destino de execução.
\end{block}
\vspace{3mm}
\begin{enumerate}
\item Extrair estimativas de volume e características estruturais da consulta.
\item Calcular um score e selecionar \texttt{small}, \texttt{medium} ou \texttt{large}.
\item Integrar análise, decisão e histórico em um protótipo reproduzível.
\item Verificar a lógica e delimitar o que ainda depende de avaliação integrada.
\end{enumerate}
\source{Objetivos do TCC e módulos do protótipo.}
\end{frame}

\begin{frame}{Referencial teórico}
\begin{description}
\item[Apache Iceberg] Snapshots, manifestos e estatísticas organizam os metadados das tabelas.
\item[Trino] O \texttt{EXPLAIN (TYPE IO)} descreve os objetos acessados e suas estimativas; estimativa não equivale a bytes físicos medidos.
\item[Trino Gateway] Provê um ponto de entrada para múltiplos clusters e políticas de roteamento.
\item[Trabalho relacionado] Strausz et al. (2025) investigam seleção de engines com modelo de custo aprendido.
\end{description}
\vspace{2mm}
\textbf{Abordagem do DyraSQL:} regra explícita de pontuação entre capacidades do mesmo motor, com pesos configuráveis.
\source{Apache Iceberg, Table Spec; Trino, EXPLAIN; Trino Gateway; arXiv:2506.02802.}
\end{frame}

\begin{frame}{Metodologia e arquitetura do protótipo}
\begin{center}
\begin{tikzpicture}[node distance=4mm and 5mm, every node/.style={font=\small,align=center},
box/.style={draw=teal,rounded corners=2pt,fill=paper,minimum height=12mm,text width=25mm},
arr/.style={-{Latex[length=2mm]},thick,draw=teal}]
\node[box] (sql) {Consulta SQL\\Cliente / Proxy};
\node[box,right=of sql] (analysis) {Análise SQL\\e EXPLAIN};
\node[box,right=of analysis] (score) {DyraSQL Core\\Score / decisão};
\node[box,right=of score] (trino) {Trino\\small / medium / large};
\node[box,below=8mm of score] (pg) {PostgreSQL\\Cache / histórico};
\node[box,below=8mm of trino] (lake) {Iceberg REST\\MinIO};
\draw[arr] (sql)--(analysis);\draw[arr] (analysis)--(score);\draw[arr] (score)--(trino);
\draw[arr] (score)--(pg);\draw[arr] (pg.west)-| (analysis.south);
\draw[arr] (trino)--(lake);
\end{tikzpicture}
\end{center}
Pesquisa aplicada com desenvolvimento e verificação experimental do protótipo.
\begin{itemize}
\item Identificação da consulta por fingerprint; consulta ao cache antes da análise.
\item Na ausência de cache: análise, cálculo dos fatores e armazenamento da decisão.
\end{itemize}
\source{Código em src/; diagrama conceitual, sem indicação de serviços ativos.}
\end{frame}

\begin{frame}{Como o algoritmo seleciona o destino}
\begin{equation*}
\boxed{S = 0{,}5 f_v + 0{,}3 f_c + 0{,}2 f_h}
\end{equation*}
\begin{columns}[T,onlytextwidth]
\column{.49\textwidth}
\begin{block}{Fatores normalizados}
\begin{itemize}
\item $f_v$: tamanho e linhas estimados.
\item $f_c$: operações e filtros SQL.
\item $f_h$: histórico de execução; valor padrão 0,5 quando ausente.
\end{itemize}
\end{block}
\column{.46\textwidth}
\begin{block}{Regra de seleção}
\begin{tabular}{ll}
\toprule Faixa & Destino\\\midrule
$S < 0{,}3$ & \texttt{small}\\[2mm]
$0{,}3 \leq S \leq 0{,}7$ & \texttt{medium}\\[2mm]
$S > 0{,}7$ & \texttt{large}\\\bottomrule
\end{tabular}
\end{block}
\end{columns}
\vspace{3mm}
Os pesos são parâmetros fixos configuráveis nesta versão; a implementação não treina um modelo de aprendizado.
\source{decision\_engine.py; history\_manager.py.}
\end{frame}

\begin{frame}{Resultados: execução local da decisão}
\evidence{02_decisoes.png}{execução do código do DecisionEngine; entradas controladas em decisoes.json.}
\end{frame}

\begin{frame}{Resultados: requisições à API e cache}
\evidence{04_api_cache.png}{TestClient da aplicação; Trino e PostgreSQL simulados nesta coleta.}
\end{frame}

\begin{frame}{Resultados: verificação automatizada}
\evidence{06_testes_linux.png}{pytest no WSL/Ubuntu, checkout isolado LF; testes\_linux\_lf.txt e linux\_lf.xml.}
\end{frame}

\begin{frame}{Limitações e alcance das evidências}
\begin{block}{O que foi demonstrado nesta coleta}
Execução do cálculo de score, classificação em três faixas, análise sintática e respostas HTTP da API com dependências controladas.
\end{block}
\vspace{2mm}
\begin{itemize}
\item A inicialização integrada foi bloqueada no download das imagens MinIO.
\item Não há nova execução distribuída de consultas nem medição de custo, memória ou latência de EXPLAIN nesta coleta.
\item Testes com mocks verificam o comportamento do código; não comprovam a integração com serviços reais.
\end{itemize}
\source{docker\_compose.txt; testes\_completos.txt; testes\_linux.txt.}
\end{frame}

\begin{frame}{Cronograma das próximas atividades}
\small
\begin{tabular}{p{24mm}p{67mm}p{39mm}}
\toprule Período & Atividade & Entregável\\\midrule
Situação atual & Artigo finalizado pelo autor; material do pitch preparado. & Slides e evidências locais.\\[3mm]
Antes de gravar & Revisar resultados e texto com o orientador; ensaiar a apresentação. & Roteiro entre 5 e 7 minutos.\\[3mm]
Gravação & Narrar os slides com identificação do autor e do orientador. & Vídeo de apresentação.\\[3mm]
Antes da entrega & Publicar como não listado; registrar URL e obter autorização prévia do orientador. & PDF do link e entrega autorizada.\\\bottomrule
\end{tabular}
\vspace{4mm}
\normalsize\color{muted}Sequência de conclusão proposta; datas dependem do prazo da disciplina e da revisão do orientador.
\end{frame}

\begin{frame}{Contribuição e encerramento}
\begin{block}{Contribuição do trabalho}
Um protótipo de roteamento com regra explícita que combina volume, estrutura SQL e histórico para selecionar entre capacidades do Trino.
\end{block}
\vspace{4mm}
\textbf{Resultado verificável:} lógica e API exercitadas localmente, com logs, entradas e respostas preservados.
\vspace{4mm}

\textbf{Continuidade técnica:} reproduzir a infraestrutura com imagens disponíveis e avaliar desempenho e custo com cargas controladas.
\vspace{5mm}

{\large\color{teal}Obrigado!}\par
\small Aristides Henrique Gonçalves da Cruz\quad|\quad Orientador: Gustavo Luís Soares
\end{frame}

\appendix
\begin{frame}{Apoio: análise sintática executada}
\evidence{03_analise_sql.png}{análise real com QueryAnalyzer e SQLGlot; consulta não executada em banco.}
\end{frame}

\begin{frame}{Apoio: execução inicial no Windows}
\evidence{01_testes.png}{saída original do pytest no Windows; falhas de portabilidade preservadas no log.}
\end{frame}

\begin{frame}{Apoio: inicialização integrada bloqueada}
\evidence{05_integracao.png}{docker compose; diagnóstico de acesso aos registros em diagnostico\_minio.txt.}
\end{frame}

\begin{frame}{Apoio: logs da aplicação local}
\evidence{07_logs.png}{execucao\_local.txt; chamadas em processo e dependências externas simuladas.}
\end{frame}

\begin{frame}{Apoio: diferenças entre ambientes de teste}
\evidence{08_portabilidade.png}{logs originais Windows, WSL com CRLF e checkout Linux LF.}
\end{frame}

\begin{frame}{Apoio: diagnóstico de acesso às imagens MinIO}
\evidence{09_registro_minio.png}{consultas anônimas reais aos registros; diagnostico\_minio.txt.}
\end{frame}

\begin{frame}{Referências principais}
\footnotesize
\begin{thebibliography}{9}
\bibitem{iceberg} APACHE SOFTWARE FOUNDATION. \textbf{Apache Iceberg Table Spec}.\par
\url{https://iceberg.apache.org/spec/}
\bibitem{trino} TRINODB. \textbf{EXPLAIN}. Documentação Trino.\par
\url{https://trino.io/docs/current/sql/explain.html}
\bibitem{gateway} TRINODB. \textbf{Trino Gateway: Overview}.\par
\url{https://trinodb.github.io/trino-gateway/}
\bibitem{strausz} STRAUSZ, A.; PARDON, N.; GIURGIU, I. \textbf{A Learned Cost Model-based Cross-engine Optimizer for SQL Workloads}. 2025.\par
\url{https://arxiv.org/abs/2506.02802}
\end{thebibliography}
\source{Fontes consultadas em 13 out. 2026. Slides de apoio fora do roteiro cronometrado.}
\end{frame}
\end{document}
